\hypertarget{absolute-basics-c98}{%
\chapter{ABSOLUTE BASICS (C++98)}\label{absolute-basics-c98}}

\hypertarget{getting-started}{%
\section{Getting Started}\label{getting-started}}

\hypertarget{what-is-c}{%
\subsection{What is C++?}\label{what-is-c}}

C++ is a statically-typed, compiled programming language that combines
low-level memory manipulation with high-level abstractions. It's the
language of choice for performance-critical applications.

\hypertarget{your-first-program-c98}{%
\subsection{Your First Program (C++98)}\label{your-first-program-c98}}

\begin{lstlisting}[language={C++}]
#include <iostream>

int main() {
    std::cout << "Hello, World!" << std::endl;
    return 0;
}
\end{lstlisting}

\textbf{Breakdown:} - \passthrough{\lstinline!\#include <iostream>!} -
Include input/output library - \passthrough{\lstinline!std::cout!} -
Standard output stream (print to console) -
\passthrough{\lstinline!std::endl!} - End line and flush buffer -
\passthrough{\lstinline!main()!} - Entry point of program -
\passthrough{\lstinline!return 0!} - Exit code (0 = success)

\textbf{Compile and run:}

\begin{lstlisting}[language=bash]
g++ -o hello hello.cpp
./hello
\end{lstlisting}

\begin{center}\rule{0.5\linewidth}{0.5pt}\end{center}

\hypertarget{basic-types-variables}{%
\section{Basic Types \& Variables}\label{basic-types-variables}}

\hypertarget{fundamental-types-c98}{%
\subsection{Fundamental Types (C++98)}\label{fundamental-types-c98}}

\begin{lstlisting}[language={C++}]
#include <iostream>
#include <limits>

int main() {
    // Integer types
    int x = 42;                  // 32-bit integer
    short y = 10;                // 16-bit integer
    long z = 1000000;            // 32 or 64-bit integer
    long long w = 9999999999;    // 64-bit integer
    
    // Floating-point types
    float f = 3.14f;             // 32-bit (4 bytes)
    double d = 3.14159265;       // 64-bit (8 bytes)
    long double ld = 3.14159265359L;  // 80+ bits
    
    // Character and boolean types
    char c = 'A';                // Single byte
    bool b = true;               // true or false
    
    // Print sizes
    std::cout << "int: " << sizeof(int) << " bytes\n";
    std::cout << "double: " << sizeof(double) << " bytes\n";
    
    // Min/max values
    std::cout << "int max: " << std::numeric_limits<int>::max() << "\n";
    std::cout << "int min: " << std::numeric_limits<int>::min() << "\n";
    
    return 0;
}
\end{lstlisting}

\hypertarget{variable-declaration-initialization}{%
\subsection{Variable Declaration \&
Initialization}\label{variable-declaration-initialization}}

\begin{lstlisting}[language={C++}]
#include <iostream>

int main() {
    // C-style initialization
    int x = 5;
    float f = 3.14f;
    
    // Multiple variables
    int a, b, c;
    
    // Uninitialized (dangerous - contains garbage)
    int uninitialized;
    std::cout << uninitialized << "\n";  // Undefined behavior!
    
    // Constants
    const int MAX_SIZE = 100;
    // MAX_SIZE = 200;  // Error: can't modify const
    
    return 0;
}
\end{lstlisting}

\hypertarget{scope-lifetime-c98}{%
\subsection{Scope \& Lifetime (C++98)}\label{scope-lifetime-c98}}

\begin{lstlisting}[language={C++}]
#include <iostream>

int global = 100;  // Global scope - lives entire program

void function() {
    int local = 5;      // Local scope - lives only in function
    {
        int nested = 10;  // Block scope - lives only in block
        std::cout << nested << "\n";
    }
    // std::cout << nested << "\n";  // Error: nested out of scope
}

int main() {
    {
        int x = 5;
    }
    // std::cout << x << "\n";  // Error: x out of scope
    
    return 0;
}
\end{lstlisting}

\hypertarget{deep-dive-the-memory-model-of-variables}{%
\subsection{Deep Dive: The Memory Model of
Variables}\label{deep-dive-the-memory-model-of-variables}}

Understanding \emph{where} your variables live is the first step to
Godhood.

\begin{enumerate}
\def\labelenumi{\arabic{enumi}.}
\tightlist
\item
  \textbf{The Stack (Automatic Storage)}:

  \begin{itemize}
  \tightlist
  \item
    \textbf{What}: Local variables (\passthrough{\lstinline!int x!}).
  \item
    \textbf{Speed}: Extremely fast (just moving a pointer).
  \item
    \textbf{Lifetime}: Scope-based (die at
    \passthrough{\lstinline!\}!}).
  \item
    \textbf{Limit}: Small (typically 1MB-8MB). Recursion depth is
    limited by this.
  \end{itemize}
\item
  \textbf{The Heap (Dynamic Storage)}:

  \begin{itemize}
  \tightlist
  \item
    \textbf{What}: \passthrough{\lstinline!new int!},
    \passthrough{\lstinline!malloc!}.
  \item
    \textbf{Speed}: Slower (allocation requires finding free block).
  \item
    \textbf{Lifetime}: Manual (until \passthrough{\lstinline!delete!} /
    \passthrough{\lstinline!free!}).
  \item
    \textbf{Limit}: RAM size (Gigabytes).
  \end{itemize}
\item
  \textbf{Static/Global (Static Storage)}:

  \begin{itemize}
  \tightlist
  \item
    \textbf{What}: Global variables, \passthrough{\lstinline!static!}
    locals.
  \item
    \textbf{Speed}: Fast access, but initialization order is tricky.
  \item
    \textbf{Lifetime}: Program start to program end.
  \end{itemize}
\item
  \textbf{Registers}:

  \begin{itemize}
  \tightlist
  \item
    \textbf{What}: CPU internal storage.
  \item
    \textbf{Speed}: Instant (0 cycles).
  \item
    \textbf{Note}: Variables are often optimized into registers, never
    touching RAM!
  \end{itemize}
\end{enumerate}

\begin{center}\rule{0.5\linewidth}{0.5pt}\end{center}

\hypertarget{operators-control-flow}{%
\section{Operators \& Control Flow}\label{operators-control-flow}}

\hypertarget{arithmetic-operators-c98}{%
\subsection{Arithmetic Operators
(C++98)}\label{arithmetic-operators-c98}}

\begin{lstlisting}[language={C++}]
#include <iostream>

int main() {
    int a = 10, b = 3;
    
    std::cout << a + b << "\n";   // 13 (addition)
    std::cout << a - b << "\n";   // 7 (subtraction)
    std::cout << a * b << "\n";   // 30 (multiplication)
    std::cout << a / b << "\n";   // 3 (integer division)
    std::cout << a % b << "\n";   // 1 (modulo)
    
    // Compound assignment
    int x = 5;
    x += 3;   // x = 8
    x -= 2;   // x = 6
    x *= 2;   // x = 12
    x /= 3;   // x = 4
    
    // Increment/Decrement
    int y = 5;
    y++;      // Post-increment: 6
    ++y;      // Pre-increment: 7
    y--;      // Post-decrement: 6
    --y;      // Pre-decrement: 5
    
    return 0;
}
\end{lstlisting}

\hypertarget{deep-dive-bitwise-mastery-low-level-optimization}{%
\subsection{Deep Dive: Bitwise Mastery (Low-Level
Optimization)}\label{deep-dive-bitwise-mastery-low-level-optimization}}

Bitwise operators manipulate individual bits. Essential for embedded
systems, graphics, and cryptography.

\hypertarget{the-operators}{%
\subsubsection{The Operators}\label{the-operators}}

\begin{itemize}
\tightlist
\item
  \passthrough{\lstinline!\&!} (AND): Both bits must be 1.
\item
  \passthrough{\lstinline!|!} (OR): At least one bit must be 1.
\item
  \passthrough{\lstinline!\^!} (XOR): Bits must be different.
\item
  \passthrough{\lstinline!\~!} (NOT): Flip all bits.
\item
  \passthrough{\lstinline!<<!} (Left Shift): Multiply by 2\^{}N.
\item
  \passthrough{\lstinline!>>!} (Right Shift): Divide by 2\^{}N.
\end{itemize}

\hypertarget{god-tier-tricks}{%
\subsubsection{God-Tier Tricks}\label{god-tier-tricks}}

\begin{enumerate}
\def\labelenumi{\arabic{enumi}.}
\tightlist
\item
  \textbf{Check Odd/Even}: \passthrough{\lstinline!(x \& 1) == 0!}
  (Even). Faster than \passthrough{\lstinline!\% 2!}.
\item
  \textbf{Multiply by 2}: \passthrough{\lstinline!x << 1!}.
\item
  \textbf{Divide by 2}: \passthrough{\lstinline!x >> 1!}.
\item
  \textbf{Clear Last Set Bit}: \passthrough{\lstinline!x \& (x - 1)!}.
  Used to count set bits (Kernighan's Algorithm).
\item
  \textbf{Check Power of 2}:
  \passthrough{\lstinline!(x > 0) \&\& ((x \& (x - 1)) == 0)!}.
\item
  \textbf{Toggle Bit N}: \passthrough{\lstinline!x \^= (1 << N)!}.
\item
  \textbf{Set Bit N}: \passthrough{\lstinline!x |= (1 << N)!}.
\item
  \textbf{Clear Bit N}: \passthrough{\lstinline!x \&= \~(1 << N)!}.
\end{enumerate}

\begin{lstlisting}[language={C++}]
// Fast Power of 2 check
bool isPowerOf2(int x) {
    return x && !(x & (x - 1));
}
\end{lstlisting}

\hypertarget{comparison-logical-operators-c98}{%
\subsection{Comparison \& Logical Operators
(C++98)}\label{comparison-logical-operators-c98}}

\begin{lstlisting}[language={C++}]
#include <iostream>

int main() {
    int a = 10, b = 5;
    
    // Comparison (return true/false)
    std::cout << (a > b) << "\n";   // 1 (true)
    std::cout << (a < b) << "\n";   // 0 (false)
    std::cout << (a == b) << "\n";  // 0 (false)
    std::cout << (a != b) << "\n";  // 1 (true)
    std::cout << (a >= b) << "\n";  // 1 (true)
    std::cout << (a <= b) << "\n";  // 0 (false)
    
    // Logical operators
    bool x = true, y = false;
    std::cout << (x && y) << "\n";  // 0 (AND)
    std::cout << (x || y) << "\n";  // 1 (OR)
    std::cout << (!x) << "\n";      // 0 (NOT)
    
    return 0;
}
\end{lstlisting}

\hypertarget{if-else-statements-c98}{%
\subsection{If-Else Statements (C++98)}\label{if-else-statements-c98}}

\begin{lstlisting}[language={C++}]
#include <iostream>

int main() {
    int score = 85;
    
    // Basic if-else
    if (score >= 90) {
        std::cout << "Grade: A\n";
    } else if (score >= 80) {
        std::cout << "Grade: B\n";
    } else if (score >= 70) {
        std::cout << "Grade: C\n";
    } else {
        std::cout << "Grade: F\n";
    }
    
    // Ternary operator
    std::string grade = (score >= 80) ? "Pass" : "Fail";
    std::cout << grade << "\n";
    
    return 0;
}
\end{lstlisting}

\hypertarget{loops-c98}{%
\subsection{Loops (C++98)}\label{loops-c98}}

\begin{lstlisting}[language={C++}]
#include <iostream>

int main() {
    // While loop
    int i = 0;
    while (i < 5) {
        std::cout << i << " ";
        i++;
    }
    std::cout << "\n";
    
    // Do-while loop (executes at least once)
    int j = 0;
    do {
        std::cout << j << " ";
        j++;
    } while (j < 5);
    std::cout << "\n";
    
    // For loop
    for (int k = 0; k < 5; k++) {
        std::cout << k << " ";
    }
    std::cout << "\n";
    
    // Break and continue
    for (int m = 0; m < 10; m++) {
        if (m == 3) continue;  // Skip 3
        if (m == 7) break;     // Exit at 7
        std::cout << m << " ";
    }
    std::cout << "\n";
    
    return 0;
}
\end{lstlisting}

\hypertarget{switch-statement-c98}{%
\subsection{Switch Statement (C++98)}\label{switch-statement-c98}}

\begin{lstlisting}[language={C++}]
#include <iostream>

int main() {
    int day = 3;
    
    switch (day) {
        case 1:
            std::cout << "Monday\n";
            break;
        case 2:
            std::cout << "Tuesday\n";
            break;
        case 3:
            std::cout << "Wednesday\n";
            break;
        default:
            std::cout << "Unknown day\n";
    }
    
    return 0;
}
\end{lstlisting}

\begin{center}\rule{0.5\linewidth}{0.5pt}\end{center}

\hypertarget{functions}{%
\section{Functions}\label{functions}}

\hypertarget{function-declaration-definition-c98}{%
\subsection{Function Declaration \& Definition
(C++98)}\label{function-declaration-definition-c98}}

\begin{lstlisting}[language={C++}]
#include <iostream>

// Function declaration (prototype)
int add(int a, int b);
void print_hello();

// Function definition
int add(int a, int b) {
    return a + b;
}

void print_hello() {
    std::cout << "Hello!\n";
}

int main() {
    print_hello();
    std::cout << add(5, 3) << "\n";  // 8
    return 0;
}
\end{lstlisting}

\hypertarget{parameters-return-values-c98}{%
\subsection{Parameters \& Return Values
(C++98)}\label{parameters-return-values-c98}}

\begin{lstlisting}[language={C++}]
#include <iostream>

// Pass by value (copy)
void increment_value(int x) {
    x++;
    std::cout << "Inside: " << x << "\n";
}

// Pass by reference (same variable)
void increment_ref(int& x) {
    x++;
    std::cout << "Inside: " << x << "\n";
}

// Pass by const reference (can't modify)
void print_value(const int& x) {
    std::cout << x << "\n";
}

// Returning by value
int get_value() {
    return 42;
}

// Returning by reference (dangerous!)
int& get_global() {
    static int x = 100;
    return x;
}

int main() {
    int a = 5;
    
    increment_value(a);   // Copy passed
    std::cout << a << "\n";  // Still 5
    
    increment_ref(a);      // Reference passed
    std::cout << a << "\n";  // Now 6
    
    print_value(a);        // Can't modify a
    
    return 0;
}
\end{lstlisting}

\hypertarget{default-parameters-c98}{%
\subsection{Default Parameters (C++98)}\label{default-parameters-c98}}

\begin{lstlisting}[language={C++}]
#include <iostream>

void greet(const std::string& name = "World") {
    std::cout << "Hello, " << name << "!\n";
}

int main() {
    greet();                // Uses default: "World"
    greet("Alice");         // Uses provided: "Alice"
    return 0;
}
\end{lstlisting}

\hypertarget{function-overloading-c98}{%
\subsection{Function Overloading
(C++98)}\label{function-overloading-c98}}

\begin{lstlisting}[language={C++}]
#include <iostream>

// Same function name, different parameters
int add(int a, int b) {
    return a + b;
}

double add(double a, double b) {
    return a + b;
}

void print(int x) {
    std::cout << "Integer: " << x << "\n";
}

void print(double x) {
    std::cout << "Double: " << x << "\n";
}

void print(const std::string& s) {
    std::cout << "String: " << s << "\n";
}

int main() {
    std::cout << add(5, 3) << "\n";      // 8 (int version)
    std::cout << add(2.5, 3.7) << "\n";  // 6.2 (double version)
    
    print(42);           // Integer version
    print(3.14);         // Double version
    print("Hello");      // String version
    
    return 0;
}
\end{lstlisting}

\begin{center}\rule{0.5\linewidth}{0.5pt}\end{center}

\hypertarget{arrays-pointers}{%
\section{Arrays \& Pointers}\label{arrays-pointers}}

\hypertarget{arrays-c98}{%
\subsection{Arrays (C++98)}\label{arrays-c98}}

\begin{lstlisting}[language={C++}]
#include <iostream>

int main() {
    // Array declaration and initialization
    int arr[5] = {1, 2, 3, 4, 5};
    
    // Access elements (0-indexed)
    std::cout << arr[0] << "\n";  // 1
    std::cout << arr[4] << "\n";  // 5
    
    // Array size (local arrays only)
    int size = 5;
    for (int i = 0; i < size; i++) {
        std::cout << arr[i] << " ";
    }
    std::cout << "\n";
    
    // 2D array
    int matrix[3][3] = {
        {1, 2, 3},
        {4, 5, 6},
        {7, 8, 9}
    };
    
    std::cout << matrix[1][2] << "\n";  // 6
    
    return 0;
}
\end{lstlisting}

\hypertarget{pointers-c98}{%
\subsection{Pointers (C++98)}\label{pointers-c98}}

\begin{lstlisting}[language={C++}]
#include <iostream>

int main() {
    int x = 42;
    
    // Create a pointer
    int* ptr = &x;  // & = address-of operator
    
    // Dereference pointer
    std::cout << *ptr << "\n";   // 42 (dereference with *)
    
    // Modify through pointer
    *ptr = 100;
    std::cout << x << "\n";      // 100
    
    // Pointer arithmetic
    int arr[5] = {10, 20, 30, 40, 50};
    int* p = arr;           // Array decays to pointer
    
    std::cout << p[0] << "\n";     // 10
    std::cout << *(p+1) << "\n";   // 20
    std::cout << p[2] << "\n";     // 30
    
    // Null pointer
    int* null_ptr = nullptr;  // or NULL in C++98
    
    // Pointer to pointer
    int** pp = &ptr;
    std::cout << **pp << "\n";  // 100
    
    return 0;
}
\end{lstlisting}

\hypertarget{dynamic-memory-c98}{%
\subsection{Dynamic Memory (C++98)}\label{dynamic-memory-c98}}

\begin{lstlisting}[language={C++}]
#include <iostream>

int main() {
    // Allocate single value on heap
    int* ptr = new int;
    *ptr = 42;
    std::cout << *ptr << "\n";
    delete ptr;          // Must deallocate
    ptr = nullptr;       // Good practice
    
    // Allocate array on heap
    int* arr = new int[10];
    for (int i = 0; i < 10; i++) {
        arr[i] = i * i;
    }
    delete[] arr;        // Note the [] for arrays
    
    // Forgetting to delete = memory leak
    int* leaked = new int(100);
    // delete leaked;  // Forgot this!
    
    return 0;
}
\end{lstlisting}
